\scanpage{40}
% archive foliation: 39

Dans le cas qui nous intéresse ($\unsure{\zeta}\in\mathrm{Br}(\mathbb Q)$ d'indice $n=2$, ramifié exactement en $\unsure{2|p}$ et $\infty$) \deleted{\unsure{on trouve}}
on trouve que $\unsure{\operatorname{Brd}}(C)$ est la \deleted{groupe} \inserted{catégorie} de Picard splittée d'invariants $\mathbb Z/2\mathbb Z$, $\mathbb Z/2\mathbb Z$, le générateur $\rho$ de $\pi_0$ \unsure{correspond} au \deleted{\unsure{l'endomorphisme}} foncteur qui à une courbe elliptique $E$ associe $E^{(p)}=E/\operatorname{Ker}F_E$)
correspond à l'unique idéal premier $\mathfrak m$ de $O$ sur $p$, et on a
\diagram{\begin{tikzcd}
N_{O/\mathbb Z}(\mathfrak m) \arrow[r] \arrow[d,phantom,"\parallel"] &
N_{O/\mathbb Z}(O)=\mathbb Z\\
p\mathbb Z &
\end{tikzcd}}
, on choisit $\bar\rho=p$.

Ceci nous dit que \uline{si $p$ est inversible dans} $A'$, alors l'image de $(\rho,\bar\rho)$ dans
\[
\unsure{\operatorname{Brd}}(C_{A'},T_{A'})
\simeq\deleted{\unsure{\operatorname{Brd}(C_{A'})}}
H^1(S',\mu_2)
\]
est la \unsure{classe} de $p$ modulo des carrés.

Par contre, prenons le cas où $A'=\mathbb Z_p$, alors
\deleted{\unsure{on est dans un groupe}}
\[
\unsure{\operatorname{Brd}}(C'_p,T'_p)
\simeq\unsure{\operatorname{Brd}}(C_{A'},T_{A'})
\]
est une extension de $\mathbb Z/2\mathbb Z$ par
$\mathbb Z_p^*/\mathbb Z_p^{*2}$
($\simeq\mathbb Z/2\mathbb Z$ si $p\ne2$), en fait c'est
\[
\underbrace{\mathbb Z_p^*/\mathbb Z_p^{*2}}_
{\parallel\ H^1(\mathbb Z_p,\mu_2)\ \deleted{\unsure{(\mathbb Z_p,A)}}}
\times
\underbrace{\mathbb Z/2\mathbb Z}_{\text{engendré par }(\mathfrak m,p)},
\]
\margin{$C_p=C'$\\$T_p=T'$}
et les classes d'isom possibles des \inserted{(du moins à l'$\infty$)}
\deleted{\unsure{liens}}
forment un torseur sur ce groupe
(\deleted{\unsure{le plus grand}}
$\simeq\mathbb Z/2\mathbb Z\times\mathbb Z/2\mathbb Z$ si $p\ne2$).

D'autre part, si $A'=\mathbb R$, on peut considérer les classes d'isomorphie de la gerbe de lien $\mu_2$ sur $\mathbb R$ définie par $C$ (soit $\unsure{\mathfrak g}_{\mathbb R}$, où $\unsure{\mathfrak g}$ est la gerbe sur $\mathbb Z[1/p]$) et l'ensemble des classes d'\inserted{\unsure{morphismes d'}}équivalences avec la gerbe \unsure{type} sur $\mathbb R$ ($=$ ensemble des \uline{\unsure{plexions|flexions}} de la gerbe), qui est un torseur sous
\[
H^1(\mathbb R,\mu_2)=\mathbb R^*/\mathbb R^{*2}.
\]
On trouve donc
